\scanpage{59}\mspage{8}
propre et lisse sur $S$ à fibres géom. connexes de dim 1, genre 1, et $\sigma$ un autom induisant $-\mathrm{id}$ sur $\underline{\mathrm{Pic}}^0_{E/S}$ \uline{équivaut} à celle \inserted{de $(E,a)$, où $a$ section de $\underline{\mathrm{Pic}}^2_{E/S}$ ($\simeq E\mathbin{\wedge}_{E_0}E$), ou encore à celle} \deleted{de} $(E_0,\underline I)$, où $E_0$ est une courbe elliptique relative sur $S$, et où $I$ est un torseur sous \deleted{$\,{}_2E_0\overset{\mathrm{déf}}{=}$} $\underline V\overset{\mathrm{déf}}{=}{}_2(E_0)$. On aura
\[
 \left\{\begin{aligned}
 &E=\underline I\mathbin{\wedge}_{\underline V}E_0\\
 &\sigma\ \text{déduit de la symétrie }x\mapsto-x\ \text{de }E_0\\
 &\hspace{2em}\text{via la formule précédente}
 \end{aligned}\right. \tag*{(21)}
\]
On aura une opération de $\underline V\deleted{\subset}={}_2E_0\subset E_0$ sur $E$, et un isom can.
\[
 \begin{aligned}
 &E/\underline V\qquad\simeq\qquad E_0\\
 &\qquad\wr\!|\\
 &E\mathbin{\wedge}_{E_0}E\simeq\underline{\mathrm{Pic}}^2_{E/S}
 \end{aligned} \tag*{(22)}
\]
% slanted
\margin{\uline{NB} L'opération de $\underline V$ sur $E$ commute à l'action de $\sigma$}
isom qui commute aux symétries $\sigma$ sur $E$, $\sigma_0:x\mapsto-x$ sur $E_0$.

Considérons les schémas quotients
\[
 X=E/\{1,\sigma\},\qquad\Sigma=E_0/\{1,\sigma_0\} \tag*{(23)}
\]
ce sont des fibrés en \deleted{\unsure{droites}} \inserted{droites} projectives, \deleted{et on a les} $V$ opère \inserted{(monomorphiquement)} sur $X$ \deleted{et sur $\Sigma$} via son \deleted{\unsure{(24)}} action sur $E$, et on aura
\[
 X/V\simeq\Sigma, \tag*{(24)}
\]
d'ailleurs le diviseur de ramification pour
\[
 E\longrightarrow X\qquad(\text{rev. quadratique}) \tag*{(25)}
\]
% archive sheet number: 58
